\section{Matroid Products}\label{sec:MatProds}

In this section we will revisit some ideas of Mason \cite{MGlueing} and Las Vergnas \cite{LV81} particularly with respect to tensor and symmetric powers of matroids, working with matroids rather than valuated matroids.
   
Work with matroid products in the category of matroids with strong maps was explored by Crapo and Rota in \cite{CR70}, and it was known at the time of these works that matroid products did not exist in this category. However, this is not the notion of matroid product that we are discussing in this paper. Instead, we interest ourselves in matroid products in a weaker sense, as defined by Las Vergnas for tensor products, and by Mason for symmetric powers.
%%%%%%%%%%%%%%%%%%%%%%%%%%%%%%%%%%%%%%%%%%%%%%%%%%%%%%%

%%%%%%%%%%%%%%%%%%%%%%%%%%%%%%%%%%%%%%%%%%%%%%%%%%%%%%%

%%%%%%%%%%%%%%%%%%%%%%%%%%%%%%%%%%%%%%%%%%%%%%%%%%%%%%%

%%%%%%%%%%%%%%%%%%%%%%%%%%%%%%%%%%%%%%%%%%%%%%%%%%%%%%%
\subsection{Tensor Products}\label{subsec:tensorProducts}
    Given matroids $M$ and $N$, we aim to construct a new matroid, called a tensor product of $M$ and $N$, which will be a matroid with ground set equal to the Cartesian product $E(M)\times E(N)$ with rank $\rk(M)\rk(N)$ and inheriting some structure from $M$ and $N$. While this operation is always possible for matroids represented over a common field $k$, it is not necessarily unique, nor is it  possible for matroids in general as shown by LasVergnas in \cite{LV81}*{Proposition 2.1}. 
    
    Several results presented in this subsection have been known since LasVergnas's original paper, however, we revisit this work and introduce a new proof of \cite{LV81}*{Proposition 2.1}.
    
    \begin{definition}\label{def:tensorproduct}
    Let $M$ and $N$ be loopless matroids on ground sets $E(M)$ and $E(N)$ respectively. A matroid $P$ on ground set $E(M)\times E(N)$ is called a \textbf{quasi product} of $M$ and $N$ if for every non-loop $e\in E(M)$ we have
    \begin{equation} \label{eqn:tensorRight}
        N\cong P|_{\set{(e,x)\mid x\in E(N)}}
    \end{equation}
    and for every non-loop $f\in E(N)$ we have
    \begin{equation}\label{eqn:tensorLeft}
        M\cong P|_{\set{(x,f)\mid x\in E(M)}},
    \end{equation}
    where these isomorphisms are explicitly given by the maps $x\mapsto (e,x)$ and $x\mapsto (x,f)$ respectively. If $e$ is a loop in $E(M)$ and $f$ a loop in $E(N)$ then we have
    \begin{equation}\label{eqn:tensorRight0}
        \rk(P|_{\set{(e,x)\mid x\in E(N)}})=0
    \end{equation}
    or
    \begin{equation}\label{eqn:tensorLeft0}
    \rk(P|_{\set{(x,f)\mid x\in E(M)}})=0
    \end{equation}
    respectively.

    We will refer to the set $\set{(e,x)\mid x\in E(N)}$ as the \textbf{row} associated to $e\in E(M)$ and we will refer to the set $\set{(x,f)\mid x\in E(M)}$ as the \textbf{column} associated to an element $f\in E(N)$ .
    \end{definition}
    
    Isomorphisms \ref{eqn:tensorRight} and \ref{eqn:tensorLeft} imply that, given bases $B_1$ of $M$ and $B_2$ of $N$, that $B_1\times B_2$ is a spanning set of any quasi product of $M$ and $N$. Therefore the rank of a quasi product is bounded above by $\rk(M)\rk(N)$, and this bound can be obtained whenever $M$ and $N$ are $F$-realizable matroids. We call that a quasi product of $M$ and $N$ with rank $\rk(M)\rk(N)$ a \textbf{tensor product} of $M$ and $N$. Moreover, if $M$ and $N$ are the same matroid, then we refer to any of their tensor products as a \textbf{second tensor power}.

    The following lemma is a result of LasVergnas, and is often useful for understanding the structure of tensor products.
    
    \begin{lemma}\label{lem:tensorCharacterizations}\cite{LV81}[Lemma 2.2]
    Let $M$ and $N$ be matroids and let $P$ be a quasi product of $M$ and $N$. The following are equivalent.
    \begin{enumerate}
        \item $\rk(P)=\rk(M)\rk(N).$
        \item For every basis $B_M$ of $M$ and every basis $B_N$ of $N$, $B_M\times B_N$ is a basis of $P$.
        \item For every pair of flats $F'\subset F$ of $M$ and $G'\subset G$ of $N$ 
        $$F'\times G \cup F\times G'$$
        is a flat of $P$.
        \item For every pair of sets $X'\subset X\subset E_M$ and $Y'\subset Y\subset E_N$
        $$\rk(X'\times Y\cup X\times Y')= \rk(X'\times Y)+\rk(X\times Y)- \rk(X'\times Y').$$
    \end{enumerate}
    \end{lemma}
    
    The motivating example for the definition of matroid tensor products is, as is often the case, linear spaces. Given two matroids $M_1$ and $M_2$ that are represented as linear subspaces of $L(M_1)\subset \mathbb{F}^{d_1}$ and $L(M_2)\subset \mathbb{F}^{d_2}$, the tensor product $L(M_1)\otimes L(M_2)$ is a linear subspace of $\mathbb{F}^{d_1d_2}$ whose corresponding matroid is a tensor product of $M_1$ and $M_2$. Given matrix representations of $L(M_1)$ and $L(M_2)$ their Kronecker product yields a matrix representation of $L(M_1)\otimes L(M_2)$.
    
    It is important to keep in mind that matroid tensor products may not be unique. In fact, even for a fixed field $\mathbb{F}$ and two $k$-representable matroids $M_1$ and $M_2$, the tensor product obtained from linear subspace representations $L(M_1)$ and $L(M_2)$ may depend on the choice of representatives $L(M_1)$ and $L(M_2)$. However, existence and uniqueness do hold in some cases.
    \begin{lemma}\label{lem:tensorSums}
     Let $A,B,$ and $M$ be any matroids and let $N=A\oplus B$. If $P$ is a tensor product of $M$ and $N$ then $P= P_A\oplus P_B$ where $P_A$ is a tensor product of $M$ and $A$ and $P_B$ is a tensor product of $M$ and $B$.
    \end{lemma}
    \begin{proof}
        Consider the matroids 
        $$P_A = P|_{E(M)\times E(A)},$$
        and
        $$P_B = P|_{E(M)\times E(A)}.$$
        It is straightforward to check that $P_A$ satisfies the definition of a quasi product between $M$ and $A$. A similar claim holds for $P_B$. Thus $P_A$ and $P_B$ are quasi products of $M$ with $A$ and $B$ respectively. Moreover, since
        $$\rk(P)=\rk(M)(\rk(A)+\rk(B))=\rk(M)\rk(A)+\rk(M)\rk(B),$$
        we may conclude that $\rk(P_A)=\rk(M)\rk(A)$ is a tensor product and $\rk(P_B)= \rk(M)\rk(B)$ is also a tensor product. 
    \end{proof}

    \begin{lemma}\label{lem:tensorSimplification}
    Let $M$ and $N$ be any matroids. Then any quasi product of $M$ with $N$ depends only the simplification of $M$. 
    \end{lemma}
    \begin{proof}
        Let $P$ be a quasi product of $M$ with $N$. Suppose $M$ has a loop $e$. Then for any $x$ in $E(N)$, we have that $(e,x)$ is a loop of $P$. This is an immediate consequence of Equation \ref{eqn:tensorRight0}. In particular, adding or removing a loop $e$ from $M$ affects the quasi product $P$ by adding or removing loops $\set{e}\times E(N)$ from $P$. 

        Now suppose that $S$ is a parallelism class of $M$. Then Isomorphism \ref{eqn:tensorRight} implies that $S\times \set{x}$ is contained in a parallelism class $S'$ of $P$. In particular, replacing the parallelism class $S$ with a representative $e$ affects $P$ by replacing the subset $S\times \set{x}$ of the parallelism class $S'$ with $(e,x)$. 
    \end{proof}
\begin{proposition}\label{prop:uniqueTensors}
Let $M$ be any matroid. Then there is a unique tnesor product between $M$ and $U_{n,n}$.
\end{proposition}
    \begin{proof}
            If $P$ is a tensor product of $M$ and $U_{n,n}$, then Proposition \ref{lem:tensorSums} implies that $P = \oplus_{i=1}^n P_i$ where $P_i$ is a tensor product of $M$ and $U_{1,1}$. In particular $P_i\cong M$ and so $P=\oplus_{i=1}^n M$. This proves existence and uniqueness. It is easy to see that $\rk(P)=n\rk(M)=\rk(U_{n,n})\rk(M)$ and so $P$ is tensor product.
    \end{proof}
\commentout{  
    \begin{proposition}\label{prop:uniqueTensors}
        Let $M$ be any matroid and let $N$ be $U_{r,n}$ where $r\in \set{0,1,n}$. Then a tensor product of $M$ and $N$ exists and is unique.
    \end{proposition}  
    \begin{proof}
    If $P$ is a quasi product of $M$ and $U_{0,n}$ then it is always the case that $\rk(P|_{\set{(e,x)\mid x\in [n]}})=0$. In particular every element is a loop and therefore $\rk(P)=0$ and so $P$ is necessarily isomorphic to the matroid $U_{0, n|E(M)|}$, which is indeed a tensor product.

    If $P$ is a quasi product of $M$ and $U_{1,n}$ then for any $e\in [n]$ it is necessarily the case that
    $$P|_{\set{(x,e)\mid x\in E(M)}} = M$$
    $$P|_{\set{(e,x)\mid x\in [n]}} = U_{1,n}.$$
    Moreover, the rank of $P$ is at most $\rk(M)$. We claim that $P$ is necessarily the matroid on $E(M)\times [n]$ with parallelism classes $\set{e\times [n]}_{e\in E(M)}$ and whose simplification is $M$, call this matroid $Q$. Note that any quasi power of $M$ and $U_{1,n}$ contains the circuit set of $Q$. In particular, if $P$ is not equal to $Q$, $P$ has strictly more circuits and $\rk(P)<rk(Q)=\rk(M)$. However, the rank of a quasi-product is bounded below by $\rk(M)+\rk(N)-1=\rk(M)$ \cite{LV81}*{Lemma 1.2}, a contradiction. Thus $Q$ exists and is unique, and it is a tensor product.

    
    If $P$ is a tensor product of $M$ and $U_{n,n}$, then Proposition \ref{lem:tensorSums} implies that $P = \oplus_{i=1}^n P_i$ where $P_i$ is a tensor product of $M$ and $U_{1,1}$. In particular $P_i\cong M$ and so $P=\oplus_{i=1}^n M$. This proves existence and uniqueness. It is easy to see that $\rk(P)=n\rk(M)=\rk(U_{n,n})\rk(M)$ and so $P$ is tensor product.

    \end{proof}
}

    \begin{corollary}\label{cor:uniqueSimplification}
        Let $M$ be any matroid and let $N$ be a matroid whose simplification is $U_{n,n}$. Then a tensor product between $M$ and $N$ exists and is unique.
    \end{corollary}
    \begin{proof}
    This follows from Propositions \ref{lem:tensorSums} and \ref{prop:uniqueTensors}, along with Lemma \ref{lem:tensorSimplification}
    \end{proof}
 
    Little is known about the existence of matroid tensor products beyond Corollary \ref{cor:uniqueSimplification} as it is seemingly easier to demonstrate when a tensor product does not exist than to construct a non-representable example.
    
    In his first work on matroid tensor products, Las Vergnas proved that the V\'amos matroid $V_8$ has no second tensor power. However, in \cite{LV81} this proof is simplified into two smaller results. First, \cite{LV81}*{Proposition 2.1} proves that no tensor product between $U_{2,3}$ and $V_8$ exists. Then \cite{LV81}*{Remark 2.3} provides that tensor products are minor-closed in the sense that if $M$ and $N$ have a tensor product $P$, then for any minors $M'$ of $M$ and $N'$ of $N$, $M'$ and $N'$ have a tensor product $P'$ that is a minor of $P$. Since $U_{2,3}$ is a minor of $V_8$ we conclude that no second tensor power of $V_8$ exists. In particular, the V\'amos matroid is a forbidden minor of the class of matroids with second tensor powers.

    We will provide a new proof of this result using a slightly more general technique while preserving the original notation. Recall that the V\'amos matroid is the rank-4 paving matroid on the $8$ element set $\set{a, a', b, b', c, c', d, d'}$, which is often abbreviated as $A\cup B\cup C\cup D$ where $X:=\set{x,x'}$, whose cyclic hyperplanes, which are the hyperplanes that are unions of circuits, are the sets $\set{A\cup B, A\cup C, A\cup D, B\cup C, B\cup D}.$ For paving matroids these cyclic hyperplanes together with their rank uniquely determine the paving matroid. We present a simplified proof that no tensor product between $U_{2,3}$ and $V_8$ exists.

    \begin{proposition}\label{prop:XandY}
        Let $[3]$ denote the ground set of $U_{2,3}$ and let $E:=A\cup B\cup C\cup D$ denote the ground set of $V_8$. Let $P$ be any quasi product of $U_{2,3}$ and $V_8$. Then either the set
        $$X:=([3]\times A)\cup (1\times B)\cup (2\times C)\cup (3\times D)$$
        or the set
        $$Y:= ([3]\times B)\cup (1\times A)\cup (2\times C)\cup (3\times D),$$
        is a spanning set.
    \end{proposition}
    \begin{proof}
    
        Let $P$ be a quasi product of  $U_{2,3}$ and $V_8$. If one of $X$ or $Y$ is a spanning set, we're done. So suppose neither $X$ nor $Y$ is spanning. We claim that $X$ and $Y$ are both hyperplanes of $P$. To see this, we show that adding any element to either $X$ or $Y$ results in a spanning set. To make this more clear we refer to the ground set in terms of rows, indexed by the ground set $[3]$ of $U_{2,3}$, and columns indexed by the ground set $E$ of $V_8$. We first consider $X$. Adding a element $e=(e_1,e_2)$ to $X$ leads to a subset of the form
            $$e_1\times (A\cup T\cup e_2),$$
    where $T\in \set{B, C, D}$. Now the closure of $X\cup e$ necessarily contains all of row $e_1$ as $A\cup T\cup e_2$ is a spanning set of $E$ in $V_8$. This now implies that $\ol{X\cup e}$ contains at least $2$ elements in each column of $[3]\times E$, which provides a spanning set of $U_{2,3}$ in each column. Thus the closure in every column is the entire column and so $\ol{X\cup e}=[3]\times E$. A similar argument holds for $Y$, where $Y\cup e$ contains the row $e_1\times (B\cup T\cup e_2)$ where $T\in{A,C,D}$. We conclude that $X$ and $Y$ are hyperplanes of $P$. 

    
    Now consider that $Z:= X\cap Y$ is also a corank-1 set. This follows as as adding the element $e = (1,c)$ leads to a spanning set in row 1. Following the previous argument we see that $\ol{Z\cup (1,c)}$ contains at least 2 elements in the columns of $C$ and $D$, which is a spanning set in those column. So $\ol{(Z\cup(1,c))}$ contains the set $[3]\times (C\cup D)$. However, $C\cup D$ is a basis of $V_8$, and so every row of $\ol{(Z\cup(1,c))}$ contains a basis. Therefore $\ol{(Z\cup(1,c))}= [3]\times E$. This is a contradiction as the intersection of two distinct hyperplanes, namely $X$ and $Y$, must have corank at least $2$, whereas $Z$ has corank $1$. Thus either $X$ or $Y$ must be a spanning set.
    \end{proof}
    
    \begin{proposition}\label{prop:VxU23}\cite{LV81}*{Proposition 2.1}
    Any quasi product of $U_{2,3}$ and $V_8$ has rank at most $7$. In particular, no tensor product between $U_{2,3}$ and $V_8$ exists.
    \end{proposition}
    \begin{proof}
        Let $[3]$ denote the ground set of $U_{2,3}$ and let $E:=A\cup B\cup C\cup D$ denote the ground set of $V_8$. Let $P$ be a quasi product of $U_{2,3}$ and $V_8$ and let the set $X$ be as in proposition \ref{prop:XandY}. Without loss of generality let $X$ be spanning. We prove $X$ has rank at most $7$. To see this, consider that $X$ is necessarily spanned by $([2]\times A)\cup \set{(1,b), (2,c), (3,d)}$, which contains $7$ elements, thus $\rk(P)=\rk(X)\leq 7$.
    \end{proof}
    
    As a consequence of Proposition \ref{prop:VxU23}, the V\'amos matroid has no tensor powers, and we further note the following small corollary
    
    \begin{corollary}\label{cor:VxU23}
    Let $V_8$ be the V\'amos matroid and $N$ be any matroid. Then a tensor product between $V_8$ and $N$ exists if and only if the simplification of $N$ is a free matroid. 
    \end{corollary}
    \begin{proof}
    If the simplification of $N$ is not free, then $N$ contains a circuit of size at least 3, and therefore $N$ contains $U_{2,3}$ as a minor and Proposition \ref{prop:VxU23} implies no tensor product between $V_8$ and $N$ exists. Otherwise, the simplification of $N$ is free and Corollary \ref{cor:uniqueSimplification} implies that $V_8$ and $N$ have a unique tensor product.
    \end{proof}
       
    While the connection between tensor products and tropical geometry is not as immediate as will be seen with symmetric powers, we are able to use tensor products in order to prove certain symmetric powers do not exist, as done in \cite{DR19}.
%%%%%%%%%%%%%%%%%%%%%%%%%%%%%%%%%%%%%%%%%%%%%%%%%%%%%%%

%%%%%%%%%%%%%%%%%%%%%%%%%%%%%%%%%%%%%%%%%%%%%%%%%%%%%%%

%%%%%%%%%%%%%%%%%%%%%%%%%%%%%%%%%%%%%%%%%%%%%%%%%%%%%%%

%%%%%%%%%%%%%%%%%%%%%%%%%%%%%%%%%%%%%%%%%%%%%%%%%%%%%%%

\subsection{Symmetric Powers}\label{subsec:SymPowers}

Matroid symmetric powers were first introduced by Mason \cite{MGlueing} where it was shown, as with tensor powers, that not every matroid has large symmetric powers. As briefly noted in the introduction, the non-existence of a tensor product for the matroids $M$ and $N$ implies the non-existence of a second symmetric power of $M\oplus N$. Mason's construction does not rely on this connection between symmetric and tensor products of matroids, showing that some connected matroids lack large symmetric powers.

Recall that if $E$ is any set, that we use the notation $\Sym_d(E)$ to denote the $d\th$ symmetric power of $E$. That is, $\Sym_d(E)$ is the set of unordered $d$-tuples of elements of $E$. Moreover, if we have an element $\alpha \in \Sym_i(E)$ we can consider the map
    \begin{align*}
        *_\alpha: \Sym_d(E)&\to \Sym_{d+i}(E)\\
        \beta&\mapsto \alpha\beta,
    \end{align*}
 which we think of as a ``multiplication map." When working with symmetric powers of subsets $S\subset \Sym_{i}(E)$ and $T\subset \Sym_{d}(E)$, we will write 
$$S\cdot T:= \set{ st\in \Sym_{d+i}(E)\mid s\in S, \spc t\in T}.$$

The most familiar and pertinent example of symmetric powers arises from commutative polynomial rings. Given a polynomial ring $k[x_1,\dots,x_n]$ over a field $k$ (or semifield), the degree $d$ monomials in $k[x_1,\dots,x_n]$ exactly correspond to the set $\Sym_d(\set{x_1,\dots,x_n}),$ which we denote by $\Mon_d(n)$ --- the degree $d$ monomials in $n$ variables. The multiplication map $*_\alpha$ is then given by multiplication by the monomial $x^\alpha$ in $k[x_1,\dots,x_n]$.

In the setting of tropical polynomial semirings, we will see matroids as collections of linear tropical polynomials, and their $d\th$ symmetric powers, as defined in Definition \ref{def:symprod}, as collections of degree-$d$ tropical polynomials.

\begin{remark}
In Definition \ref{def:symprod} we define a $d\th$ symmetric power of a fixed valuated matroid $\nu$ to be a sequence $(\nu_1,\dots,\nu_d)$ satisfying a a few conditions. However, Isomorphism \ref{eqn:symprodEmbedding} provides that $\nu_d$ uniquely determines the entire sequence $(\nu_1,\dots,\nu_d)$, so nothing is lost when referring to $\nu_d$ alone. However, $\nu$ may have multiple $d\th$ symmetric powers, as we will see in Example \ref{ex:matrixSymprodExample}, and when constructing a $d\th$ symmetric power it will be useful to explicitly record the entire sequence $(\nu_1,\dots,\nu_d)$. We will further discuss the space of all symmetric powers of a fixed valuated matroid $\nu$ in Section \ref{subsec:SpacesofSymPow}.
\end{remark}

\begin{remark}
If $(\nu_1,\dots,\nu_d)$ is a $d\th$ symmetric power of $\nu$, then for every $1\leq i\leq d$ we have that the truncation $(\nu_1,\dots,\nu_i)$ is an $i\th$ symmetric power of $\nu$.
\end{remark}



\begin{remark}\label{rem:loopless}
Isomorphism \ref{eqn:symprodEmbedding} implies that if $\nu$ is loop-free, then so is every symmetric quasi power. Moreover, Equation \ref{eqn:symprodLoops} implies that adding or removing loops to $\nu$ to obtain a matroid $\nu'$ uniquely determines how to add or remove loops from any given $d\th$ symmetric quasi power $(\nu_1,\dots,\nu_d)$ to obtain a $d\th$ symmetric quasi power of $\nu'$. In particular, we may assume without loss of generality that $\nu$ is loop-free and all subsequent results extend to matroids with loops.
\end{remark}

It is important to note that we refer to \textit{a} symmetric power of a matroid rather than \textit{the} symmetric power, as it may not be unique, not even for realizable matroids.

The following is stated in \cite{MGlueing}, but we restate it here and include a brief proof for intuition's sake.

\begin{proposition}\cite{MGlueing}\label{prop:masonFullRank}
    For a matroid $M$ of rank $r$ the rank of any symmetric quasi power $M_d$ is at most $\binom{r+d-1}{d}$.
\end{proposition}

\begin{proof}
For any symmetric quasi power $M_d$ of $M$ and any basis $B$ of $M$, Isomorphism \ref{eqn:symprodEmbedding} implies that the set $\Sym_d(B)$ is a spanning set of $M_d$. Thus the rank of $M_d$ can never be greater than $|\Sym_d(B)|=\binom{r+d-1}{d}$. 
\end{proof}
 In light of Proposition \ref{prop:masonFullRank}, we see that symmetric powers are exactly those symmetric quasi-powers that achieve the upper bound on the rank. Moreover given any field $k$, every $k$-representable matroid $M$ has a $k$-representable $d\th$ symmetric power that can be directly constructed. Given any $k$-matrix $L$ representing $M$ we can construct the $d\th$ \textbf{Macaulay matrix} of $L$, which will be a representation of a $d\th$ symmetric power of $M$. We outline the construction of the second Macaulay matrix here. Let $r\leq n$ and let $L$ be a rank-$r$ $r\times n$ matrix over a field $k$. The second Macaulay matrix of $L$, which we will denote $\Sym_2(L)$, is an $\binom{r+1}{2}\times \binom{n+1}{2}$ matrix whose rows are indexed by $\Sym_2([r])$ and whose columns are indexed by $\Sym_2([n])$. Fixing the standard basis $e_1,\dots,e_r$ of $k^r$, we can write each column $C_j$ of $L$ as a formal sum 
$$C_j = \sum_{i=1}^r L_{i,j}e_i.$$
We can now consider $C_j$ as a linear form in the polynomial ring $k[u_1,\dots,u_r]$ under the inclusion $\phi:k^n\mapsto k[u_1,\dots,u_r]$, $e_i\mapsto u_i$.

The column $ij$ of $\Sym_2(L)$ is computed by taking the product $\phi(C_i)\phi(C_j)\in k[u_1,\dots,u_r]_2$, which is then a formal sum in the basis $\Mon_2(r)$, i.e in basis vectors indexed by $\Sym_2([r])$. In particular we express 
$$\phi(C_i)\phi(C_j)=\sum_{x,y\in[r]} c_{xy,ij}u_xu_y.$$

We now define the entries of the second Macaulay matrix $\Sym_2(L)$ to be, for $x,y\in [r]$ and $i,j\in[n]$
$$\Sym_2(L)_{xy,ij} = c_{xy,ij},$$
that is, the $xy$ coefficient of the product of the polynomials associated to columns $C_i$ and $C_j$. This definition extends to defining $\Sym_{d}(L)$ by considering products in $k[x_1,\dots,x_n]_d.$

As indicated in our notation above, the $d\th$ Macaulay matrix of $L$ is indeed a matrix that represents a $d\th$ symmetric power of the matroid $M$ that $L$ represents. We will demonstrate this for the second Macaulay matrix. 
\begin{example}\label{ex:matrixSymprodExample}
Consider the matrix representing $U_{2,4}$ over any field other than $\F_2$ with $\alpha \neq 0,1$
$$M:=\begin{bmatrix}
1 & 0 & 1 & 1\\
0 & 1 & 1 & \alpha
\end{bmatrix}.$$
The second Macaulay matrix of $M$ is a matrix whose column space represents a second symmetric power of $M$. It's columns are indexed by pairs of columns of $M$ and its rows are indexed by pairs of column entries of some fixed basis of $M$, in this case we use the first two columns of $M$.
$$\Sym_2(M):=\kbordermatrix{
        & c_{11}& c_{12}& c_{13}& c_{14} & c_{22}& c_{23}& c_{24} & c_{33}& c_{34}  & c_{44}  \\
r_{11}  &1      & 0     & 1     & 1      & 0     & 0     & 0      & 1     & 1       & 1       \\
r_{12}  &0      & 1     & 1     & \alpha & 0     & 1     & 1      & 2     & \alpha+1& 2\alpha \\
r_{22}  &0      & 0     & 0     & 0      & 1     & 1     & \alpha & 1     & \alpha       & \alpha^2}.$$
Note that the symmetric power of the fixed basis $\set{c_1,c_2}$ is  $\set{c_{11},c_{12},c_{22}}$, which is a basis of $\Sym_2(M),$ so it is indeed rank $3$ as desired.

Consider the minor $B$ consisting of the following columns of $\Sym_2(M)$
$$B:=\kbordermatrix{
        &c_{11} & c_{22}    & c_{34}    \\
r_{11}  &1      & 0         & 0         \\
r_{12}  &0      & 0         & \alpha +1    \\
r_{22}  &0      & 1         & \alpha  }.$$
We can see that $B$ has determinant $-(\alpha+1)$. In particular if the characteristic is not $2$ then we may choose $\alpha= -1$ and $B$ will not be a basis. Otherwise, if we can choose $\alpha\neq 1,0,-1$ then $B$ will be a basis. Similarly, the matrix
$$C:=\kbordermatrix{
        &c_{11} & c_{22}    & c_{44}    \\
r_{11}  &1      & 0         & 1         \\
r_{12}  &0      & 0         & 2\alpha    \\
r_{22}  &0      & 1         & \alpha^2  }.$$
has determinant $-2\alpha$, which is 0 if and only if $k$ has characteristic 2. In particular is never a basis when represented in characteristic $2$ and it is always a basis otherwise.
\end{example}

\begin{remark}
Example \ref{ex:matrixSymprodExample} shows that the underlying matroid of the second Macaulay matrix depends both on the field over which we take the representation $L$, as well as the representation $L$ itself over any fixed field. However, if $L$ is a representation of $M$ over the field $F$, then the second Macaualay matrix is preserved under action of $\GL_r(F)$ by left-multiplication in the following sense. If $U\in \GL_r(F)$ we have
    $$\Sym_2(UL) = \Sym_2(U)\Sym_2(L).$$

Example \ref{ex:matrixSymprodExample} demonstrates that the analogue for right-hand multiplication by $V\in \GL_n(F)$ does not hold.

\end{remark}

\subsection{Symmetric Powers by Flats}\label{subsec:sympowByFlats}
In this section we restrict our view to matroids rather than valuated matroids, as we will be predominantly talking about flats of matroids, which are not yet axiomatized for valuated matroids.

In Definition \ref{def:symprod} we decided on an axiomatization for matroid symmetric powers that is potentially stronger than that which is proposed by Mason. Originally motivated from Lov\'asz's early work with symmetric powers of realizable matroids \cite{Lovasz77}, Mason proposed the following condition for a symmetric quasi power, as defined in Definition \ref{def:symprod}, to be a symmetric power.
\begin{equation}\label{axiom:flatAxiom} 
\text{For every flat } F \text{ of } M \text{ the set } F\cdot \Sym_{d-1}(E) \text{ is a flat of } M_d.
\end{equation}

In Definition \ref{def:symprod} we instead required that Equation \ref{eqn:symprodRank} be satisfied.

\begin{proposition}\label{prop:SymprodImpliesFlats}
    If $M$ is a valuated matroid and $(M_1,\dots, M_d)$ is a $d\th$ symmetric power in the sense of Definition \ref{def:symprod}, then $M_d$ satisfies Axiom \ref{axiom:flatAxiom}.
\end{proposition}
\commentout{
\begin{proof}
Suppose that $M_d$ is a a symmetric power of $M$. Let $H$ be a hyperplane of $M$ and suppose that $H^d:=H\cdot \Sym_{d-1}(E)$ is not a flat of $M_d$. Then there is an $e=\prod_{i=1}^d e_i\in \ol{H^d}\sm H_d$ and $e_i\notin H$ for all $1\leq i\leq d$. We claim that $e_1^d\in \ol{H^d}$. Consider that $\ol{H^d}|_{e/e_2} \supset \ol{H\cup e_2}=E$. In particular, we conclude that $e_1^2\prod_{i=3}^d e_i\in \ol{H^d}$. Let
$$f_j = \prod_{i=1}^j e_1 \prod_{i={j+1}}^d e_i$$
we may repeat this argument for the restriction $\ol{H^d}|_{f_j/e_{j+1}}$ to conclude that $f_{j+1}\in \ol{H^d}$. In particular $f_d = e_1^d\in \ol{H^d}.$

Now let $I$ be a basis of $H$. Then $B:=I\cup e_1$ is a basis of $M$. This implies that $\Sym_d(B)$ is a basis of $M_d$, or in particular that $I\cdot \Sym_{d-1}(B) = \Sym_d(B)-e_1^d$ is a basis of $H^d$. In particular, $e_1^d$ is not in $\ol{H^d}$, a contradiction. So $H^d$ must be a flat.


We now note that every flat $F$ is an intersection of hyperplanes: 
$$F=\bigcap_{i=1} H_i$$
and that 
$$F\cdot \Sym_{d-1}(E) = \bigcap_{i=1} H_i\cdot \Sym_{d-1}(E)$$
is an intersection of flats, and therefore a flat of $M_d$, as desired.
\end{proof}
}

We will prove Proposition \ref{prop:SymprodImpliesFlats} in Section \ref{subsec:sympowCharacterizations}, using our new geometric techniques to avoid an in-depth matroid theoretical argument. While we are able to prove that Definition \ref{def:symprod} implies Axiom \ref{axiom:flatAxiom}, the reverse is not at all obvious. Despite substantial effort, we were not able to make progress on this problem and have left it as an open problem left by Lov\'asz and Mason.

\begin{conjecture}(Lovasz '77, Mason '81)\label{conj:flatcondition}
If $(M_1,\dots,M_d)$ is a symmetric quasi power of $M$ and for every $F\in \sF(M)$ we have $F\cdot \Sym_{d-1}(E)\in \sF(M_d)$, then $M_d$ has rank $\binom{\rk(M)+d-1}{d}$.
\end{conjecture}

In order to prove the conjecture, it is sufficient to prove the following auxiliary result.

\begin{conjecture}\label{conj:flatcondition2}
    If $M$ is a matroid on the ground set $E$, $(M_1,\dots,M_d)$ is a symmetric quasi power of $M$, and for every $F\in \sF(M)$ we have $F\cdot \Sym_{d-1}(E)\in \sF(M_d)$, then there is an $e\in E$ such that the contraction
    $$M_d/(e\cdot \Sym_{d-1}(E))$$
    is a symmetric quasi power of $M/e$.
\end{conjecture}

Conjecture \ref{conj:flatcondition2} is nontrivial in general, as it is possible for $(M_1,\dots,M_d)$ to be a symetric quasi power, that is not full rank, and for the contraction $M_d/(e\cdot \Sym_{d-1}(E))$ to not be a symmetric quasi power for some $M/e$. However, it is not clear if this is necessarily possible that this could be the case for all $e$. It is not clear if Conjecture \ref{conj:flatcondition} is true, not even for second symmetric powers,  without additional substantial hypotheses.


\subsection{Further characterizations of Matroid Symmetric Powers}\label{subsec:sympowCharacterizations}

In this section we will discuss further characterisations of symmetric powers. In particular, we will provide a geometric characterization, alongside a characterisation by flats that requires a strictly stronger condition than that of Conjecture \ref{conj:flatcondition}.

In order to establish a geometric interpretation of matroid symmetric powers, we will need to define the \textit{variety} of a symmetric power, which follows a similar construction as the Bergman fan of the matroid, but is meant to instead capture the relationship of $M_d$ to $M$ rather than the distinct geometry of $M_d$

\begin{definition}\label{def:matroidVariety}
Let $M$ be a matroid on ground set $[n]\cong\set{x_1,\dots,x_n}$, and let $M_d$ be a symmetric quasi power of $M$ on ground set $\Sym_{d}([n])\cong \Mon_d(n)$. For each circuit $C$ of $M_d$ we define the \textit{circuit polynomial}
    $$p_C := \bigoplus_{x\in C} x\in \T[x_1,\dots,x_n]_d.$$

We can now define the \textbf{variety} of $M_d$ as
$$V(M_d) : = \bigcap_{C\in \sC(M_d)} V\left(p_C\right)\subset \T^n.$$
\end{definition}
The primary difference between the Bergman fan of $M_d$ and its variety is that the Bergman fan treats each element of $\Sym_d(E)$ as an independent variable whereas the variety takes into account the algebraic dependencies between the variables. Thus $B(M_d)\subset\T^{|\Sym_d(E)|}$ and $V(M_d)\subset \T^n$. Moreover, we notice that it is always the case that $V(M_d)\subset B(M)$ and that $V(M_d)$ maps to $B(M_d)$ through the Veronese embedding.


\begin{proposition}\label{prop:symprodCharacterisations}
Let $M_d$ be a $d\th$ symmetric quasi power of $M$. The following are equivalent for $M_d$
\begin{enumerate}
    \item $M_d$ has rank $\binom{\rk(M)+d-1}{d}$
    \item For every basis $B$ of $M$, $\Sym_d(B)$ is a basis of $M_d$
    \item The variety $V(M_d)$ is supported on $B(M)$.
    \item For every $\omega\in B(M)$ the set
    $$\set{x\in \Sym_d(E)\mid x(\omega)>0}$$
    is a flat of $M_d$
\end{enumerate}
\end{proposition}
The equivalence between $1)$ and $2)$ follows from Proposition \ref{prop:masonFullRank}. So we need only prove points $3)$ and $4)$. We will prove $3)$ in Section \ref{sec:LinTropIdeals} and prove $4)$ in this section.

As mentioned before, Proposition \ref{prop:symprodCharacterisations}\textcolor{red}{.4} is stronger than Axiom \ref{axiom:flatAxiom}, which can be phrased in the following way.
\begin{quote}
    For every $\omega\in B(M)\cap \R^n_{\geq 0}$. The set
    $$\set{x\in \Sym_d(E)\mid x(\omega)>0}$$
    is a flat of $M_d$.
\end{quote}
Before proceeding with the proof of \ref{prop:symprodCharacterisations}\textcolor{red}{.4} recall that the lineality space of $B(M)$ contains $\R \mathbbm{1}$. In particular, the positive part of the Bergman fan determines the entire fan. The discrepancy between Conjecture \ref{conj:flatcondition} and Proposition \ref{prop:symprodCharacterisations}\textcolor{red}{.4} is simply a matter of whether or not we need to consider the flats corresponding to every point in the Bergman fan, or just the positive part --- the part that determines the entire fan.

\begin{proof} We prove that 1) is equivalent to 3) in the Section \ref{sec:LinTropIdeals}. Thus we focus on the equivalence between 3) and 4).

Suppose that $V(M_d)\neq B(M)$. Then $V(M_d)\subsetneq B(M)$ and so we let $\omega\in B(M)\sm V(M_d)$. Since $\omega\notin V(M_d)$ there is a circuit $C$ of $M_d$ that attains its minimum exactly once at $\omega$, say at the monomial $x$. Thus $x(\omega)$ is the unique minimum value of $C(\omega)$. Since $B(M)$ has lineality space $\R\mathbbm{1}$ we may translate $\omega$ to $\omega' = \omega- \frac{x(\omega)}{d} \mathbbm{1}$ so that $x(\omega')=0$ is the unique minimum value of $C$ at $\omega'$. Now consider that this implies the set
$$S:=\set{y\in \Sym_d(E)\mid y(\omega')>0}$$
contains all of $C$ except for $x$. Thus $S$ is not a flat, as desired.

Now suppose that the set 
$$S = \set{y\in \Sym_d(E)\mid y(\omega)>0}$$
is a not a flat for some $\omega\in B(M)$. We claim that $V(M_d)\neq B(M)$. Since $S$ is not a flat there is a circuit $C$ such that $|C\sm S|=\set{x}$ --- a singleton. In particular, $x(\omega)$ is the unique minimum of $C$ at $\omega$ and so $\omega\notin V(M_d)$ and $V(M_d)\neq B(M)$, as desired.
\end{proof}

\begin{proof}[Proof of Proposition \ref{prop:SymprodImpliesFlats}]
If $M$ is a matroid and $(M_1,\dots,M_d)$ is a symmetric power of $M$ as in Definition \ref{def:symprod}, then by Proposition \ref{prop:symprodCharacterisations}, for every $\omega\in B(M)$ we have that

$$S_\omega:= \set{x\in \Sym_d(E): x(\omega)>0}.$$
Let $F$ be a flat of $M$. The set $F\cdot \Sym_{d-1}(E)$ is the set of elements of $\Sym_d(E)$ divisible by some element of $F$. Now let $\omega= e_F$. A monomial $x$ satisfies $x(\omega)>0$ if and only if $x$ is divisible by an element of $F$. In particular $F\cdot\Sym_{d-1}(E)=S_\omega$ and is therefore a flat of $M_d$, as desired.
\end{proof}

\subsection{Open Questions}\label{subsec:openQuestions}

We conclude this section by posing a few open questions that we were not successful in pursing in this paper. In the case of linear spaces there is a well-established connection between tensor and symmetric powers of linear spaces via the quotient of the action of the symmetric groups, we pose this as a question in the context of matroid powers

\begin{question}
Is there a relationship, e.g via matroid quotients, between matroid tensor powers and matroid symmetric powers? Are there suitable hypothesis on a matroid tensor power to guarantee the existence of a matroid symmetric power?
\end{question}

 We might further refine the previous question and ask if symmetric powers are, in some sense,  easier to construct than tensor powers.
\begin{question}\label{question:symprodsWithoutTensor}
Can a matroid have a $d \th$ symmetric power and no $d\th$ tensor power, or vice-versa?
\end{question}

A prime candidate for answering this question is to determine if the V\'amos matroid $V_8$ has a second symmetric power. We imagine that it does not, but if it does in fact have one, it would resolve this question and also prove that the class of matroids with second symmetric powers is not closed under duality. As a quick proof of this, if $V_8$ has a second symmetric power $P$, then so would $V_8\oplus U_{1,3}$ via $P\oplus (V_8\otimes U_{1,3})\oplus (U_{1,3})_2$, where $V_8\otimes U_{1,3}\cong V_8$. However, since $V_8$ is self-dual, the dual of $V_8\oplus U_{1,3}$ is $V_8\oplus U_{2,3}$, whose second symmetric power does not exist \cite{DR19}*{Theorem 5.2}. We refine this into its own question.

\begin{question}\label{question:vamosSym2}
Does the V\'amos matroid have a second symmetric power?
\end{question}
%%%%%%%%%%%%%%%%%%%%%%%%%%%%%%%%%%%%%%%%%%%%%%%%%%%%

%%%%%%%%%%%%%%%%%%%%%%%%%%%%%%%%%%%%%%%%%%%%%%%%%%%%

%%%%%%%%%%%%%%%%%%%%%%%%%%%%%%%%%%%%%%%%%%%%%%%%%%%%



%omitted:

\commentout{
It is important to note that the flat condition is false if we do not fix $M_{d-i}$ as it embeds into $M_{d}$. We demonstrate this in the following example

\begin{example}
Assume the definition of symmetric quasi power does not require all restrictions of $M_d$ to $\omega \cdot \Sym_{d-i}(E)$ to be equal for $\omega\in \Sym_i(E)$. Without this restriction, the only criteria to verify is that for every $\omega\in \Sym_{d-1}(E)$ that $M_d|_{\omega\cdot E}=M$. In particular, the matroid $P$ on $\Sym_6([4])$ whose only circuit is 
$$C:= x_1x_2x_3x_4(x_1^2+x_2^2+x_3^2+x_4^2)$$

is a symmetric quasi power of the free matroid $U_{4,4}$, as its only circuit is not divisible by any degree $5$ monomial. Hence all relevant restrictions of $P$ are isomorphic to $U_{4,4}$. This implies that all embeddings are equal to some symmetric quasi power of $U_{4,4}$, which need not be equal to each other or full rank. Now consider that since $C$ is divisible by $\prod_{i\in E}x_i$ that $C\in \cap_{F\in \sF} F\cdot \Sym_5(E)$. In particular, every $F\cdot \Sym_5(E)$ is a flat of $P$, but $P$ is not full rank.
\end{example}
}